\section{Planning Before the Moment of Action}

\begin{proposition}[Resolved entry lowers ambiguity]
If a plan provides a resolved entry specification for the target state $s_1$,
then the target ambiguity term $U(s_1)$ is weakly lower than under an otherwise
similar unplanned intention.
\end{proposition}

\begin{proof}
By definition, a resolved entry specification fixes a visible first action that
would otherwise remain undecided at the moment of action. Since $U(s_1)$ was
defined as the unresolved specification burden of entering $s_1$, resolving at
least one such burden weakly lowers $U(s_1)$.
\end{proof}

\begin{proposition}[Planning lowers future barrier under the monotone scaffold]
Suppose a plan for $s_0 \to s_1$ both lowers target ambiguity and increases
relevant preparation. Under the monotone planning assumption, the future
activation barrier is weakly lower than without that plan.
\end{proposition}

\begin{proof}
From the previous proposition, resolved entry lowers $U(s_1)$. By definition
of planning, relevant preparation is also increased. The monotone planning
assumption states that $\DV(s_0 \to s_1)$ is weakly increasing in ambiguity and
weakly decreasing in preparation. Therefore lowering ambiguity and increasing
preparation weakly lowers the activation barrier.
\end{proof}

\begin{corollary}[Last-minute self-persuasion is structurally weaker]
If last-minute self-persuasion changes neither ambiguity nor preparation, while
planning changes at least one of them in the favourable direction, then the
planned transition has weakly lower activation barrier than the merely
persuaded transition.
\end{corollary}

\begin{proof}
The merely persuaded transition leaves the relevant barrier variables unchanged
by assumption. The planned transition lowers at least one favourable term.
Apply the previous proposition.
\end{proof}

\begin{remark}
This corollary is deliberately weaker than the chapter's earlier ``planning
theorem.'' It does not claim that planning always wins in every human setting.
It states only that if planning actually changes the variables that the model
says matter, then the model predicts a lower barrier than verbal endorsement
alone.
\end{remark}
